\section{Elektr maydon potensiali}



\begin{enumerate}

    \item Bir-biridan 0,03 m masofada joylashgan va bir jinsli elektr maydonning bitta kuch chizig’ida yotgan nuqtalar orasidagi potensiallar farqi 12 V ga teng. Xuddi o’sha kuch chizig’ida bir-biridan 15 cm masofada yotgan nuqtalar orasidagi potensiallar farqini toping.
    \begin{enumerate}[label=\Alph*)]
        \item 30
        \item 120
        \item 60 % To'g'ri
        \item 15
    \end{enumerate}

    \item Yassi kondensatordagi elektr maydon kuchlanganligi 30 kV/m. Qoplamalar orasidagi potensiallar farqi 300 V. Kondensator qoplamalari orasidagi masofa (mm da) qanday?
    \begin{enumerate}[label=\Alph*)]
        \item 10 % To'g'ri
        \item 1
        \item 100
        \item 30
    \end{enumerate}

    \item Vakuumda bir-biridan 0,1 m masofada joylashgan ikkita parallel metall plastina 1 kV potensiallar farqigacha zaryadlangan. Plastinalar orasiga joylashtirilgan 100 mkC zaryadga qanday kuch ta’sir qiladi?
    \begin{enumerate}[label=\Alph*)]
        \item 0,1
        \item 1 % To'g'ri
        \item 10
        \item 5
    \end{enumerate}

    \item Yassi kondensatorning gorizontal plastinalari orasida $4,8 \cdot 10^{-12}$ kg massali chang zarrachasi muvozanatda turibdi. Agar kondensatorda kuchlanish 3000 V, plastinalar orasidagi masofa esa 2 cm bo’lsa, chang zarrachasining zaryadi elektron zaryadidan necha marta katta? ($e=1,6 \cdot 10^{-19}$ C, $g=10 \, m/s^2$).
    \begin{enumerate}[label=\Alph*)]
        \item 1000
        \item 4000
        \item 500
        \item 2000 % To'g'ri
    \end{enumerate}

    \item Yassi kondensatorning gorizontal plastinalari orasida plastmass prujinaga zaryadlangan sharcha osilgan. Kondensator 500 V kuchlanish manbaiga ulanganida, prujina yana 1 cm ga cho’ziladi. Agar prujinaning bikirligi 10 N/m, kondensator plastinalari orasidagi masofa 20 cm bo’lsa, sharchaning zaryadini (mkC da) toping.
    \begin{enumerate}[label=\Alph*)]
        \item 40 % To'g'ri
        \item 20
        \item 80
        \item 10
    \end{enumerate}

    \item Manfiy zaryadlangan $10^{-9}$ g massali chang zarrachasi yassi kondensator ichida muvozanatda turibdi. Kondensatorga 500 V kuchlanish qo'yilgan. Zarrachadan 500 ta elektron chiqib ketgandan so’ng ham u muvozanatda qolishi uchun potensiallar farqini necha voltga o’zgartirish kerak? ($d=5$ mm).
    \begin{enumerate}[label=\Alph*)]
        \item 1000
        \item 2000 % To'g'ri
        \item 500
        \item 2500
    \end{enumerate}

    \item Yassi kondensator plastinalari orasida, pastki plastinadan 10 cm masofada zaryadlangan sharcha muallaq turibdi ($U=400$ V). Agar potensiallar farqi oniy tarzda 200 V gacha kamaytirilsa, qancha vaqtdan so’ng (ms da) sharcha pastki plastinaga tushadi?
    \begin{enumerate}[label=\Alph*)]
        \item 200 % To'g'ri
        \item 100
        \item 400
        \item 50
    \end{enumerate}

    \item Elektron yassi kondensator maydoniga kuch chiziqlari yo’nalishida kirib, 0,003 m yo’lda tezligini yo'qotdi. Agar elektronning boshlang’ich tezligi va kondensator kuchlanishi 3 marta kamaytirilsa, u qanday masofada (mm da) to'xtaydi?
    \begin{enumerate}[label=\Alph*)]
        \item 0,5
        \item 3
        \item 1 % To'g'ri
        \item 0,3
    \end{enumerate}

    \item 5 kV potensiallar farqini o’tib tezlik olgan elektronlar yassi kondensatorga o’rtadan (parallel) kiradi. Elektronlar kondensatordan chiqib ketmasligi uchun qanday eng kichik kuchlanish kerak? ($L=5$ cm, $d=1$ cm).
    \begin{enumerate}[label=\Alph*)]
        \item 200
        \item 400 % To'g'ri
        \item 800
        \item 100
    \end{enumerate}

    \item Tomoni 30 cm bo’lgan teng tomonli uchburchakning ikki uchida har biri 50 nC dan bo’lgan zaryadlar bor. Uchinchi uchdagi potensialni (kV da) toping.
    \begin{enumerate}[label=\Alph*)]
        \item 1,5
        \item 6
        \item 9
        \item 3 % To'g'ri
    \end{enumerate}

    \item Tomoni 30 cm bo’lgan kvadratning ikki qarama-qarshi uchlarida 200 nC dan zaryadlar bor. Kvadratning qolgan ikki uchidagi potensialni (kV da) toping.
    \begin{enumerate}[label=\Alph*)]
        \item 12 % To'g'ri
        \item 6
        \item 18
        \item 24
    \end{enumerate}

    \item To’g’ri burchakli uchburchakning uchlarida 1, 2 va 3 nC zaryadlar bor. Agar gipotenuzaning uzunligi 20 cm bo’lsa, uning o’rtasida potensial qanchaga teng bo’ladi? ($k=9 \cdot 10^9$).
    \begin{enumerate}[label=\Alph*)]
        \item 270
        \item 540 % To'g'ri
        \item 1080
        \item 360
    \end{enumerate}

    \item Qirrasi 30 cm bo’lgan muntazam tetraedrning uchta uchida 3, 5 va –2 nC zaryadlar bor. To’rtinchi uchdagi potensialni toping.
    \begin{enumerate}[label=\Alph*)]
        \item 90
        \item 360
        \item 180 % To'g'ri
        \item 120
    \end{enumerate}

    \item Tomoni 27 cm bo’lgan muntazam oltiburchakning uchta uchida 1 nC dan, qolgan uchtasida 2 nC dan zaryadlar bor. Oltiburchakning markazidagi potensialni toping.
    \begin{enumerate}[label=\Alph*)]
        \item 150
        \item 300 % To'g'ri
        \item 600
        \item 100
    \end{enumerate}

    \item 6 cm radiusli yupqa halqa bo’ylab 4 nC zaryad taqsimlangan. Halqa o’qida, markazidan 8 cm masofada yotgan nuqtada maydon potensialini toping.
    \begin{enumerate}[label=\Alph*)]
        \item 360 % To'g'ri
        \item 180
        \item 720
        \item 450
    \end{enumerate}

    \item 30 cm radiusli sferaning sirti bo’ylab 4 nC zaryad taqsimlangan. Sfera markazida potensial qanchaga teng?
    \begin{enumerate}[label=\Alph*)]
        \item 60
        \item 240
        \item 180
        \item 120 % To'g'ri
    \end{enumerate}

    \item 10 nC musbat zaryadli sferaning markazida –5 nC manfiy zaryadli sharcha bor. Sferadan tashqarida, markazidan 9 cm masofada bo’lgan nuqtadagi potensialni toping.
    \begin{enumerate}[label=\Alph*)]
        \item 10
        \item 5 % To'g'ri
        \item 2,5
        \item 1
    \end{enumerate}

    \item Maydonning bir nuqtasidan ikkinchisiga 2 nC zaryadni ko’chirishda maydon qanday ish (mkJ da) bajaradi? 
    Nuqtalar orasidagi potesniallar farqi $\Delta \varphi = 500$ V.
    \begin{enumerate}[label=\Alph*)]
        \item 2
        \item 0,5
        \item 1 % To'g'ri
        \item 10
    \end{enumerate}

    \item Maydonning 20 V potensialli nuqtasidan 12 V potensialli nuqtasiga 8 mkC zaryadni ko’chirishda qanday ish (mkJ da) bajariladi?
    \begin{enumerate}[label=\Alph*)]
        \item 32
        \item 64 % To'g'ri
        \item 128
        \item 96
    \end{enumerate}

    \item 130 nC zaryadni cheksizlikdan maydonning qandaydur nuqtasiga ko’chirishda bajarilgan ish 65 mkJ ga teng. Shu nuqtaning potensialini(V) toping.
    \begin{enumerate}[label=\Alph*)]
        \item 250
        \item 1000
        \item 500 % To'g'ri
        \item 130
    \end{enumerate}

    \item 10 nC zaryadni cheksizlikdan zaryadlangan shar yuzasidan 20 cm masofadagi nuqtaga ko’chirishda 0,5 mkJ ish bajarildi. ($R=4$ cm). Shar sirtidagi potensialni toping.
    \begin{enumerate}[label=\Alph*)]
        \item 150
        \item 300 % To'g'ri
        \item 600
        \item 50
    \end{enumerate}

    \item Manfiy zaryadni qo’zg’almas musbat zaryad tomon ko'chirishda maydon 9 J ish bajardi va masofa yarmiga kamaydi. Shu yo’lning birinchi yarmida qanday ish (J) bajarilgan?
    \begin{enumerate}[label=\Alph*)]
        \item 3 % To'g'ri
        \item 1
        \item 4,5
        \item 6
    \end{enumerate}

\end{enumerate}